% Vietnamese translation for OI Public Library
% Source-File: IOI中国国家候选队论文/2006/汤泽/汤泽.ppt
\documentclass[11pt]{article}
\usepackage{oipl-vn}

\newcommand{\Oh}{\mathcal{O}}

\title{Tối ưu bằng tính đơn điệu\\{\large Ghi chú theo slide}}
\author{Tang Ze}
\date{2006}

\begin{document}
\maketitle

\origtitle{Monotonicity-based algorithm optimization}
\sourcepath{IOI China candidate team papers / 2006 / Tang Ze / Tang Ze.ppt}

\section{Phạm vi}

Bộ trình chiếu đi kèm bài viết về việc dùng hàng đợi hai đầu và tính đơn điệu
để tối ưu thuật toán. Ghi chú này dựa trên bản bài viết đã dịch.

\section{Quan hệ đơn điệu}

Trong nhiều công thức quy hoạch động, quyết định tối ưu của trạng thái sau
không đi lùi so với quyết định tối ưu của trạng thái trước. Khi chứng minh
được quan hệ đó, ta không cần thử lại toàn bộ các quyết định cũ cho mỗi trạng
thái.

\section{Hàng đợi ứng viên}

Ta duy trì một hàng đợi các quyết định còn có khả năng tối ưu. Đầu hàng đợi
là quyết định tốt nhất cho trạng thái hiện tại. Khi sang trạng thái mới, nếu
quyết định đầu không còn tốt hơn quyết định kế tiếp, nó bị loại. Khi thêm
quyết định mới, các quyết định ở cuối bị loại nếu chúng bị quyết định mới lấn
át trên mọi trạng thái tương lai.

\section{Ví dụ quy hoạch động}

Bài \emph{Saw Mill} minh họa dạng
\[
  f_i=\min_{j<i}\{A_j+B(j,i)\}.
\]
Tính trực tiếp tốn \(\Oh(n^2)\). Nếu điểm chuyển tối ưu đơn điệu theo \(i\),
hàng đợi ứng viên cho phép mỗi quyết định vào và ra hàng đợi số lần hằng, từ
đó giảm đáng kể thời gian.

\section{Ngoài quy hoạch động}

Cùng tư tưởng cũng xuất hiện trong cửa sổ trượt, duy trì giá trị nhỏ nhất hoặc
lớn nhất trên một đoạn đang dịch chuyển. Điều kiện để loại bỏ một phần tử là:
khi có phần tử mới tốt hơn và tồn tại lâu hơn trong các cửa sổ tương lai, phần
tử cũ không thể trở thành đáp án nữa.

\end{document}
